\documentclass[12pt]{article}
\usepackage[a4paper,margin=1in]{geometry}
\usepackage[utf8]{inputenc}
\usepackage[T1]{fontenc}
\usepackage{amsmath,amssymb}
\usepackage{textcomp}
\usepackage{parskip}
\begin{document}
\begin{center}
{\Large \textbf{2002 STEP III}}
\end{center}

\subsection*{Page 1}
\begin{flushleft}
STEP III, 2002 2\par
\bigskip\textbf{Section A: Pure Mathematics}\par
ln x\par
\bigskip\textbf{Question 1.} Find the area of the region between the curve y = and the x-axis, for 1 6 x 6 a. What\par
x\par
happens to this area as a tends to infinity?\par
ln x\par
\[Find the volume of the solid obtained when the region between the curve y = and the\]
x\par
x-axis, for 1 6 x 6 a, is rotated through 2\pi radians about the x-axis. What happens to this\par
volume as a tends to infinity?\par
a + b\par
\bigskip\textbf{Question 2.} Prove that arctan a + arctan b = arctan when 0 < a < 1 and 0 < b < 1 .\par
\lambda 1 ab \mu\par
-\par
\[Prove by induction that, for n > 1 ,\]
n\par
1 n\par
\[arctan = arctan\]
\[\lambda r^2 + r + 1 \mu \lambda n + 2 \mu\]
\[Xr=1\]
and hence find\par
\infty 1\par
arctan .\par
\[\lambda r^2 + r + 1 \mu\]
\[Xr=1\]
Hence prove that\par
\[\infty 1 \pi\]
\[arctan = .\]
\[\lambda r^2 r + 1 \mu 2\]
\[Xr=1 -\]
3 Let\par
\[f (x) = a\sqrt\{x\} \sqrt\{x\} b ,\]
- -\par
where x > b > 0 and a > 1 . Sketch the graph of f (x) . Hence show that the equation f (x) = c,\par
2 2 2\par
\[where c > 0, has no solution when c < b a 1 . Find conditions on c in terms of a and b\]
-\par
for the equation to have exactly one or exactly tw o solutions.\par
\[Solve the equations (i) 3\sqrt\{x\} \sqrt\{x\} 2 = 4 and (ii) 3\sqrt\{x\} \sqrt\{x\} 3 = 5 .\]
- - - -\par
\end{flushleft}
\subsection*{Page 2}
\begin{flushleft}
STEP III, 2002 3\par
\bigskip\textbf{Question 4.} Show that if x and y are positive and x 3 + x 2 = y 3 y 2 then x < y .\par
-\par
\[Show further that if 0 < x 6 y 1, then x 3 + x 2 < y 3 y 2 .\]
- -\par
Prove that there does not exist a pair of positive integers such that the difference of their\par
cubes is equal to the sum of their squares.\par
Find all the pairs of integers such that the difference of their cubes is equal to the sum of their\par
squares.\par
5 Give a condition that must be satisfied by p, q and r for it to be possible to write the quadratic\par
2 2\par
polynomial px + qx + r in the form p (x + h) , for some h.\par
Obtain an equation, which you need not simplify, that must be satisfied by t if it is possible\par
to write\par
2 1 2 4 3 2\par
x + bx + t x + bx + cx + dx + e\par
-\par
in the form k (x + h) , for some k and h.\par
4 3 2\par
Hence, or otherwise, write x + 6x + 9x 2x 7 as a product of two quadratic factors.\par
- -\par
\bigskip\textbf{Question 6.} Find all the solution curves of the differential equation\par
4 dy 2 2\par
\[y = y 1\]
\lambda dx \mu -\par
that pass through either of the points\par
\[\sqrt\{3\}\]
\textbf{(i)} 0, ,\par
\gamma \delta\par
\[\sqrt\{5\}\]
\textbf{(ii)} 0, .\par
\gamma \delta\par
Show also that y = 1 and y = 1 are solutions of the differential equation. Sketch all these\par
-\par
solution curves on a single set of axes.\par
\end{flushleft}
\subsection*{Page 3}
\begin{flushleft}
STEP III, 2002 4\par
\bigskip\textbf{Question 7.} Given that α and β are acute angles, show that α + β = \pi/2 if and only if cos 2 α + cos 2 β = 1.\par
In the x-y plane, the point A has coordinates (0, s) and the point C has coordinates (s, 0),\par
\[where s > 0. The point B lies in the first quadrant (x > 0, y > 0). The lengths of AB, OB\]
and CB are respectively a, b and c.\par
Show that\par
2 2 2 2 2 2 2 2 2 2\par
\[(s + b a ) + (s + b c ) = 4s b\]
- -\par
and hence that\par
2 2 2 2 2 2 2 2 2 2\par
\[(2s a c ) + (2b a c ) = 4a c .\]
- - - -\par
Deduce that\par
(a c) 2 6 2b 2 6 (a + c) 2 .\par
-\par
8 Four complex numbers u1, u2, u3 and u4 have unit modulus, and arguments θ1, θ2, θ3 and θ4,\par
\[respectively, with \pi < θ1 < θ2 < θ3 < θ4 < \pi.\]
-\par
Show that\par
\[arg (u^1 u^2) = (θ1 + θ2 \pi) + 2n\pi\]
- 2 -\par
\[where n = 0 or 1 . Deduce that\]
\[arg ((u^1 u^2) (u^4 u^3)) = arg ((u^1 u^4) (u^3 u^2)) + 2n\pi\]
- - - -\par
for some integer n.\par
Prove that\par
\[(u^1 u^2) (u^4 u^3) + (u^1 u^4) (u^3 u^2) = (u^1 u^3) (u^4 u^2) .\]
| - - | | - - | | - - |\par
\end{flushleft}
\subsection*{Page 4}
\begin{flushleft}
STEP III, 2002 5\par
\bigskip\textbf{Section B: Mechanics}\par
\bigskip\textbf{Question 9.} A tall container made of light material of negligible thickness has the form of a prism, with\par
2 3\par
a square base of area a . It contains a volume ka of fluid of uniform density. The container\par
is held so that it stands on a rough plane, which is inclined at angle θ to the horizontal, with\par
\[two of the edges of the base of the container horizontal. In the case k > tan θ, show that\]
the centre of mass of the fluid is at a distance x from the lower side of the container and at a\par
distance y from the base of the container, where\par
x 1 tan θ y k tan θ\par
\[= , = + .\]
a 2 - 12k a 2 24k\par
\[Determine the corresponding coordinates in the case k < tan θ.\]
\[The container is now released. Given that k < , show that the container will topple if\]
\[θ > 45 .\]
\textasciicircum{}\circ\par
\bigskip\textbf{Question 10.} A light hollow cylinder of radius a can rotate freely about its axis of symmetry, which is fixed\par
and horizontal. A particle of mass m is fixed to the cylinder, and a second particle, also of\par
mass m, moves on the rough inside surface of the cylinder. Initially, the cylinder is at rest,\par
with the fixed particle on the same horizontal level as its axis and the second particle at rest\par
vertically below this axis. The system is then released. Show that, if θ is the angle through\par
which the cylinder has rotated, then\par
g\par
\[θ¨ = (cos θ sin θ) ,\]
2a -\par
provided that the second particle does not slip.\par
\[3 + \sqrt\{3\}\]
Given that the coefficient of friction is , show that the second particle starts to slip\par
when the cylinder has rotated through 60 .\par
\textasciicircum{}\circ\par
\bigskip\textbf{Question 11.} A particle moves on a smooth triangular horizontal surface AOB with angle AOB = 30 . The\par
\textasciicircum{}\circ\par
surface is bounded by two vertical walls OA and OB and the coefficient of restitution between\par
the particle and the walls is e, where e < 1. The particle, which is initially at point P on the\par
\[surface and moving with velocity u^1, strikes the wall OA at M^1, with angle P M1A = θ, and\]
\[rebounds, with velocity v^1, to strike the wall OB at N^1, with angle M1N1B = θ. Find e and\]
\[v^1\]
in terms of θ.\par
\[u^1\]
The motion continues, with the particle striking side OA at M2, M3, . . . and striking side OB\par
\[at N^2, N^3, . . .. Show that, if θ < 60\]
\textasciicircum{}\circ\par
, the particle reaches O in a finite time.\par
\end{flushleft}
\subsection*{Page 5}
\begin{flushleft}
STEP III, 2002 6\par
\bigskip\textbf{Section C: Probability and Statistics}\par
\bigskip\textbf{Question 12.} In a game, a player tosses a biased coin repeatedly until two successive tails occur, when the\par
game terminates. For each head which occurs the player wins \pounds 1. If E is the expected number\par
of tosses of the coin in the course of a game, and p is the probability of a head, explain why\par
\[E = p (1 + E) + (1 p) p (2 + E) + 2 (1 p) ,\]
- -\par
and hence determine E in terms of p. Find also, in terms of p, the expected winnings in the\par
course of a game.\par
A second game is played, with the same rules, except that the player continues to toss the\par
coin until r successive tails occur. Show that the expected number of tosses in the course of\par
r\par
1 q\par
\[a game is given by the expression - , where q = 1 p.\]
pqr\par
-\par
\bigskip\textbf{Question 13.} A continuous random variable is said to have an exponential distribution with parameter λ\par
\[if its density function is f (t) = λe\]
-\par
\[λt (0 6 t < ) . If X^1 and X^2, which are independent\]
\infty\par
random variables, have exponential distributions with parameters λ1 and λ2 respectively, find\par
an expression for the probability that either X1 or X2 (or both) is less than x. Prove that if\par
\[X is the random variable whose value is the lesser of the values of X^1 and X^2, then X also\]
has an exponential distribution.\par
Route A and Route B buses run from my house to my college. The time between buses on each\par
route has an exponential distribution and the mean time between buses is 15 minutes for Route\par
A and 30 minutes for Route B. The timings of the buses on the two routes are independent.\par
If I emerge from my house one day to see a Route A bus and a Route B bus just leaving\par
the stop, show that the median wait for the next bus to my college will be approximately 7\par
minutes.\par
\end{flushleft}
\subsection*{Page 6}
\begin{flushleft}
STEP III, 2002 7\par
\bigskip\textbf{Question 14.} Prove that, for any two discrete random variables X and Y ,\par
\[Var [X + Y ] = Var [X] + Var [Y ] + 2 Cov [X, Y ] ,\]
where Var [X] is the variance of X and Cov [X, Y ] is the covariance of X and Y .\par
When a Grandmaster plays a sequence of m games of chess, she is, independently, equally\par
likely to win, lose or draw each game. If the values of the random variables W , L and D are the\par
numbers of her wins, losses and draws respectively, justify briefly the following claims:\par
\textbf{(i)} W + L + D has variance 0 ;\par
\textbf{(ii)} W + L has a binomial distribution.\par
Cov [W, L]\par
Find the value of .\par
Var [W ] Var [L]\par
p\par
\end{flushleft}
\end{document}
